% MEASUREMENT_MAP.tex, a companion to skeleton_ClassMarkers.tex of Paper A1. Written 15
% October 2026 for Paper A, and restricted on the same day to the rows that
% stay within one generation (TODO.md, pass 15).
%
% PURPOSE
%   A reference note, not part of the paper. For every object of the model it
%   records one plain-English name, its mathematical type, what it stands for
%   in the world, and whether and how an observer could obtain it. The Name
%   column is the terminology. The manuscript uses these names and no
%   synonyms, and checks/verify_measurement_map.py (suite 4) checks that every
%   term of the manuscript's Terminology appendix is a name here.
%
% TABLES
%   Inputs are the primitives the model takes as given. Internal objects are
%   built from the inputs inside the model and named for the proofs. Outputs
%   are what the rows predict.
%
% DIRECTLY OBSERVED
%   Yes exactly when the last column begins with "Recorded", and No
%   otherwise. checks/verify_measurement_map.py (MM-7) enforces the match.
%
% THE LAST COLUMN
%   Each entry begins with one of six words or phrases. "Recorded" means an
%   observer writes it down. "Elicited" means it is a belief, obtained by
%   asking the observer. "Computed" means it is calculated from recorded or
%   elicited objects. "Assumed" means the model fixes its form and no
%   measurement is proposed. "Never observed" means it exists in the world
%   but no record reaches it. "A device" means it exists only inside a proof.
%
% VERIFICATION NOTE
%   This note adds no mathematics of its own. Everything restates a model row
%   or the predecessor manuscript, and the rows are the evidence.
%
% NAMES OF RESULTS
%   Results are named by the model rows of skeleton_ClassMarkers.tex. Labels in brackets
%   such as (prior) and (threshold) are equation labels of the predecessor
%   manuscript, retired at pass 6 and readable from git at commit 91a74cb1.

\documentclass[11pt]{article}
\usepackage[T1]{fontenc}
\usepackage[utf8]{inputenc}
\usepackage{mathptmx}
\usepackage{microtype}
\usepackage[a4paper,margin=2cm]{geometry}
\usepackage{amsmath,amssymb}
\usepackage{booktabs}
\usepackage{array}
\usepackage{longtable}

\newcolumntype{P}[1]{>{\raggedright\arraybackslash}p{#1}}
% \mm{symbol}{name}{type}{represents}{directly observed}{obtained how}
\newcommand{\mm}[6]{#1 & \textbf{#2} & #3 & #4 & #5 & #6 \\ \addlinespace}
\newcommand{\lean}[1]{\texttt{\footnotesize #1}}

\title{Class Read from Group Markers\\ Measurement Map}
\author{Companion note to the manuscript}
\date{15 October 2026}

\begin{document}
\maketitle

\noindent An observer meets a person, reads a marker, and forms a belief
about the person's class through the group the marker signals (M1 to M3,
M17). Members of a group may pass, leaving the visible group (M14), and a
platform admits or excludes the group's members (M15). The discrete objects
are the group, the marker, and in the Lean rows the class. The continuous
objects are the standing, the beliefs, the refusal share and the platform's
margin. Every belief is the observer's and is a construct attributed to the
observer to rationalise what the observer does. Every rate is a fact of the
world that records can reach. The model's one payoff
is the engagement utility of M18, which enters every row only through the
engagement threshold, so the baseline value and the outside
option are read only up to the threshold they give. The sensitivity is
derived in M19 from the budget slope and the two curvatures of the payoff
over the activity level and private consumption. The engagement decision
of the predecessor manuscript, which turns a belief into a refusal, is row
M18, and the discrimination-rate function $F$ that M14 uses is derived
there. Objects that no row carries are listed with that said, so that their
names are fixed before the rows reach them.

The objects fall into three tables. Inputs are what the model takes as given,
internal objects are built from the inputs inside the model, and outputs are
what the rows predict. The column Directly observed says Yes when a record of
the world gives the object without any computation, which is when its last
column begins with Recorded, and No otherwise. An elicited belief is a No,
since the observer's answer stands in for the belief rather than showing it.
Every symbol in the manuscript's table of assumptions has exactly one row
here.

\section*{Inputs}

\footnotesize
\begin{longtable}{@{}P{1.8cm}P{2.1cm}P{2.2cm}P{3.7cm}P{1.3cm}P{3.7cm}@{}}
\caption{Inputs of the model.}\\
\toprule
\textbf{Symbol} & \textbf{Name} & \textbf{Type} & \textbf{Represents} & \textbf{Directly observed} & \textbf{Obtained how} \\
\midrule
\endfirsthead
\toprule
\textbf{Symbol} & \textbf{Name} & \textbf{Type} & \textbf{Represents} & \textbf{Directly observed} & \textbf{Obtained how} \\
\midrule
\endhead
\bottomrule
\endfoot
\mm{$g$, $\mathcal{G}$}{Group}{An element of a finite set.}{The category the
  observer reads a person into from the marker. The predecessor's host group
  and outsider group are two of them.}
  {Yes}{Recorded by the observer as the
  category the marker signals, or self-reported in a survey. In audit and
  correspondence studies it is assigned by the experimenter through the
  marker shown.}
\mm{$O$}{Outsider group}{One of the groups.}{The group the marker signals,
  whose members may pass (M14) and whom a platform admits or excludes
  (M15).}
  {Yes}{Recorded as the group the marker is read as signalling.}
\mm{$I$}{Host group}{One of the groups.}{The group alongside the outsider,
  whose clients the platform serves (M15).}
  {Yes}{Recorded as the other group in the record.}
\mm{$c$, $\mathcal{C}$}{Class}{A finite set in the Lean rows. The predecessor
  uses $[0,1]$.}{A person's economic position, income, wealth or occupation,
  as a rank.}
  {Yes}{Recorded as an income or wealth rank scaled to $[0,1]$, after
  the fact. Class is never seen at the encounter, which is the premise of the
  model.}
\mm{$y$}{Own class}{A value of class.}{The class of the observer who decides
  whether to engage (M18).}
  {Yes}{Recorded as the observer's income or wealth rank, after the fact.}
\mm{$m$}{Marker}{A signal read at the encounter.}{A trait visible at contact,
  phenotype, accent, dress or name.}
  {Yes}{Recorded as present or absent, or by its
  value. In correspondence studies it is assigned by the experimenter.}
\mm{$\pi_g(c)$}{Prior}{A probability mass function on $\mathcal{C}$ for each
  $g$. The predecessor gives it a Beta law with mean $\bar c$.}{The observer's
  belief about the class of a member of group $g$ before the marker is
  read.}
  {No}{Elicited from the observer before the marker is shown. The
  predecessor identifies it with the group's class history, its equation
  (prior), under which it is computed from the recorded classes of past
  members. The two readings can differ, and no row takes a side.}
\mm{$\ell_g(c)$}{Likelihood}{A function of $c$ and $g$ with values in
  $[0,1]$.}{The probability that the marker is seen on a person of class $c$ in
  group $g$.}
  {No}{Computed from records of the marker by group and by class, with
  class recorded after the fact. Group sufficiency restricts it to depend on
  $g$ alone (M1).}
\mm{$q(g)$}{Group weights}{A probability mass function on $\mathcal{G}$.}{The
  observer's probability that the person met belongs to group $g$, before the
  marker is read.}
  {No}{Elicited from the observer, or proxied by the shares of the
  groups in the population of the place of the encounter.}
\mm{$s_g$}{Salience}{A real number in $[0,1]$.}{The probability the observer
  assigns to the outsider group after reading the marker, the predecessor's
  $P'(g_O\mid m)$. By M3 it is a Bayesian posterior on group. Immutable markers
  put it near $1$ and mutable ones near $0$.}
  {No}{Elicited from the observer after
  the marker is shown. In audits it is the share of encounters in which the
  person carrying the marker is treated as the outsider group.}
\mm{$P$}{Passers}{A subset of a finite set of people $V$.}{The members of
  the outsider group who suppress their markers and leave the visible group
  (M14).}
  {Yes}{Recorded as the members whose marker is no longer read, where a
  panel follows people, or elicited as self-reported suppression.}
\mm{$V$}{Group members}{A finite set of people, each with a class
  $c_i$.}{The members of the outsider group in one generation, before any of
  them pass (M14).}
  {Yes}{Recorded as the people a panel or a census assigns to the
  group.}
\mm{$\mathrm{gain}$}{Passing gain}{A function of class.}{What a member of a
  given class gains by suppressing the marker and leaving the visible group
  (M14).}
  {No}{Never observed. No row derives it, and M14 uses only that it does
  not fall in class.}
\mm{$\mathrm{cost}$}{Passing cost}{A function of class.}{What suppressing
  the marker costs a member of a given class (M14).}
  {No}{Never observed. No row
  derives it, and M14 uses only that it does not rise in class.}
\mm{$W$}{Willingness to pay}{A real number.}{What an admitted outsider
  member pays the platform (M15).}
  {Yes}{Recorded as the price paid, or elicited
  as a reservation price.}
\mm{$N$}{Host clients}{A positive integer.}{The number of host-group clients
  on the platform (M15).}
  {Yes}{Recorded by the platform.}
\mm{$\alpha$}{Objective weight}{A real number in $(0,1)$.}{The weight a
  client puts on objective quality against social quality (M15).}
  {No}{Elicited
  from clients, or never observed where only admissions are seen.}
\mm{$|A|$}{Room size}{A positive integer.}{The number of attendees
  (M15).}
  {Yes}{Recorded by the platform.}
\mm{$v$}{Baseline value}{A real number.}{The engagement utility at zero
  class distance (M18).}
  {No}{Never observed. It enters every row only through the engagement
  threshold.}
\mm{$\kappa$}{Sensitivity}{A real number with $\kappa>0$,
  $\kappa=a^2\beta^2/(8(a+h))$.}{The fall in engagement utility per unit of
  squared class distance (M18), the loss each side of a pair bears at the
  compromise level (M19).}
  {No}{Computed from the private curvature, the activity curvature and the
  budget slope (M19, \lean{Budget.engagement\_utility}). It enters M18 only
  through the engagement threshold.}
\mm{$b$}{Budget}{A real number, $b=\beta c$.}{The consumption a person's
  class affords, split between the shared activity and private consumption
  (M19).}
  {Yes}{Recorded as the person's income or total spending.}
\mm{$\beta$}{Budget slope}{A real number with $\beta>0$.}{The budget per
  unit of class (M19).}
  {No}{Computed as the slope of recorded budgets on recorded class ranks.}
\mm{$p$}{Activity level}{A real number.}{The level of the one activity a
  pair shares, paid by each from the own budget (M19).}
  {Yes}{Recorded as the price per head of the shared activity.}
\mm{$\gamma$}{Activity value}{A real number.}{The payoff per unit of
  activity level (M19).}
  {No}{Never observed. It enters the rows through the ideal level and the
  affordability bound.}
\mm{$h$}{Activity curvature}{A real number with $h>0$.}{The fall in the
  activity's marginal payoff per unit of level (M19).}
  {No}{Never observed. It enters the rows through the ideal level and the
  sensitivity.}
\mm{$a$}{Private curvature}{A real number with $a>0$.}{The fall in the
  marginal payoff of private consumption per unit of consumption (M19). With
  a linear private utility the class gap costs nothing.}
  {No}{Never observed. It enters the rows through the ideal level and the
  sensitivity, and it is the curvature that a pooled budget would hide.}
\mm{$\bar u$}{Outside option}{A real number.}{The utility of declining to
  engage (M18).}
  {No}{Never observed. It enters every row only through the engagement
  threshold.}
\mm{$J$}{Observers}{A finite set of people, each with an own class.}{The
  observers over whom the refusal share is taken (M18).}
  {Yes}{Recorded as the people whose encounters with the group are
  observed.}
\mm{$\bar d$}{Engagement threshold}{A real number with $\bar d>0$.}{The
  largest expected distance at which the observer still engages, the
  predecessor's equation (threshold).}
  {No}{Elicited as the largest expected distance at which the observer
  still engages. It equals the baseline value less the outside option,
  divided by the sensitivity (M18, \lean{Engagement.engages\_iff}).}
\end{longtable}
\normalsize

\section*{Internal objects}

\footnotesize
\begin{longtable}{@{}P{1.8cm}P{2.1cm}P{2.2cm}P{3.7cm}P{1.3cm}P{3.7cm}@{}}
\caption{Objects built inside the model.}\\
\toprule
\textbf{Symbol} & \textbf{Name} & \textbf{Type} & \textbf{Represents} & \textbf{Directly observed} & \textbf{Obtained how} \\
\midrule
\endfirsthead
\toprule
\textbf{Symbol} & \textbf{Name} & \textbf{Type} & \textbf{Represents} & \textbf{Directly observed} & \textbf{Obtained how} \\
\midrule
\endhead
\bottomrule
\endfoot
\mm{$Z_g$}{Marginal}{A real number, $Z_g=\sum_c\ell_g(c)\pi_g(c)$.}{The
  probability of the marker given the group alone.}
  {No}{Computed from the
  likelihood and the prior. The marginal is also recorded directly as the frequency of
  the marker within the group.}
\mm{$Z$}{Marginal across groups}{A real number, $Z=\sum_g Z_g\,q(g)$.}{The
  probability of the marker before the group is known.}
  {No}{Computed from the
  marginals and the group weights.}
\mm{$\ell_g\equiv Z_g$}{Group sufficiency}{A restriction on the
  likelihood.}{The marker says nothing about class beyond what the group
  says. M1 shows it equals $\pi'_g=\pi_g$ under a prior with full
  support.}
  {No}{Assumed. Group sufficiency could be tested on records of the marker by class
  within one group.}
\mm{$\lambda_g$}{Constant likelihood}{A real number for each $g$.}{The
  one value the likelihood takes on every class of group $g$ when group
  sufficiency holds (M1, M2).}
  {No}{A device of M1 and M2. When the prior sums to
  one it equals the marginal $Z_g$, as M1's proof shows.}
\mm{$w$, $\xi$}{Target weights}{A probability mass function on
  $\mathcal{G}$ and a scale $\xi\neq0$.}{In M3, weights on the groups,
  and the scale of the likelihood that reaches them by Bayes' rule.}
  {No}{A device of M3, which shows any such weights are a Bayesian posterior
  on group.}
\mm{$\sum_g\pi_g(c)\,w(g)$}{Mixture}{A probability mass function on
  $\mathcal{C}$.}{The class distribution reached by weighting each group's
  prior by group weights $w$. M2 and M3 show it is the Bayesian posterior on
  class exactly when $w=q'$.}
  {No}{Computed from the priors and the weights.}
\mm{$p^*$}{Ideal level}{A real number, $p^*(b)=(ab+\gamma-1)/(a+h)$.}{The
  activity level a person would set alone, which does not fall in the
  budget (M19).}
  {No}{Computed from the budget, the activity value and the two curvatures
  (M19, \lean{Budget.ideal\_max}, \lean{Budget.ideal\_mono}). Recorded
  spending on activities against recorded budgets would test the
  monotonicity.}
\mm{$\bar p$}{Compromise level}{A real number, the midpoint of the two
  ideal levels.}{The level a pair shares, the minimiser of the sum of the two
  losses (M19).}
  {No}{Computed from the two ideal levels (M19, \lean{Budget.mid\_min}).
  The recorded price of what a pair does together against the two
  members' budgets would test it.}
\mm{$\bar c$}{Standing}{A real number in $[0,1]$.}{The mean class of a group
  in a generation. The state of the dynamic part.}
  {Yes}{Recorded as the group's
  mean class rank in that generation.}
\mm{$\nu_g$}{Spread}{A real number, $\nu_g=\sum_c\pi_g(c)(c-\bar c_g)^2$.}{The
  variance of class within group $g$ under its prior (M18).}
  {Yes}{Recorded as the variance of the recorded class ranks of the group's
  members.}
\mm{$E$}{Expected distance}{A real number,
  $E=\sum_g q'(g)\,[(\bar c_g-y)^2+\nu_g]$.}{The expected squared class
  distance from the observer to the person, under the posterior on group
  (M18).}
  {No}{Computed from the posterior on group, the standings, the spreads and
  the own class (M18, \lean{Engagement.expDist\_decomp}).}
\mm{$F$}{Discrimination-rate function}{A function $[0,1]\to[0,1]$,
  $F(\bar c)=|\{j\in J:\bar d<(\bar c-y_j)^2+\nu_O\}|/|J|$. Row M14 uses only
  that it does not rise, and M18 gives the range on which it does not.}{The
  share of the observers who refuse when the person is read as the outsider
  group with certainty, as a function of the group's standing.}
  {No}{Computed from the observers' own classes, the outsider's standing and
  spread, and the engagement threshold (M18,
  \lean{Engagement.refusal\_iff}, \lean{Engagement.refusalShare\_antitone}).
  Recorded refusal rates, callbacks, viewings or service against the group's
  recorded standing would test the two ranges.}
\mm{$V\setminus P$}{Visible group}{A finite set of people.}{The members
  of the outsider group whose markers are still read (M14).}
  {Yes}{Recorded as the
  complement of the passers.}
\mm{$\mathrm{marg}$}{Admission margin}{A real number,
  $W+N(1-\alpha)(\bar c_O-\bar c_I)/|A|$.}{The profit from admitting one more
  outsider member, the predecessor's first-order condition (M15).}
  {No}{Computed
  from its five inputs.}
\mm{$\Pi$}{Platform profit}{A function of the admission rate,
  $\Pi(\sigma)=\sigma\,\mathrm{marg}$.}{The platform's objective in the
  predecessor's first-order condition, linear in the admission rate
  (M15).}
  {No}{Computed from the admission margin. Recorded profit would test the
  linear form.}
\end{longtable}
\normalsize

\section*{Outputs}

\footnotesize
\begin{longtable}{@{}P{1.8cm}P{2.1cm}P{2.2cm}P{3.7cm}P{1.3cm}P{3.7cm}@{}}
\caption{Outputs of the model.}\\
\toprule
\textbf{Symbol} & \textbf{Name} & \textbf{Type} & \textbf{Represents} & \textbf{Directly observed} & \textbf{Obtained how} \\
\midrule
\endfirsthead
\toprule
\textbf{Symbol} & \textbf{Name} & \textbf{Type} & \textbf{Represents} & \textbf{Directly observed} & \textbf{Obtained how} \\
\midrule
\endhead
\bottomrule
\endfoot
\mm{$\pi'_g(c)$}{Posterior}{A probability mass function on
  $\mathcal{C}$.}{The observer's belief about the class of a member of group
  $g$ after the marker is read.}
  {No}{Elicited after the marker is shown. M1
  predicts that it equals the prior under group sufficiency, so a marker that
  moves it within a group refutes sufficiency for that marker.}
\mm{$q'(g)$}{Posterior on group}{A probability mass function on
  $\mathcal{G}$.}{The observer's group weights after the marker is read. For
  the outsider group it is the salience.}
  {No}{Elicited after the marker is shown,
  as for the salience.}
\mm{$P(c\mid m)$}{Posterior on class}{A probability mass function on
  $\mathcal{C}$.}{The observer's belief about the person's class after the
  marker, across groups.}
  {No}{Elicited after the marker is shown. M2 predicts the
  mixture of the priors by the posterior on group.}
\mm{$\bar c(V\setminus P)$}{Visible mean}{A real number in $[0,1]$.}{The
  mean class of the members who stay visible. M14 predicts it falls when the
  leavers are the upper members.}
  {Yes}{Recorded as the mean class rank of the
  members whose marker is still read, and compared with the mean of the
  whole group.}
\mm{$\sigma$}{Admission rate}{A real number in $[0,1]$.}{The share of
  outsider members the platform admits (M15). M15 predicts $0$ when the
  outsider stands below the admission threshold.}
  {Yes}{Recorded as the share of
  outsider applicants admitted.}
\mm{$c_E$}{Admission threshold}{A real number,
  $\bar c_I-W|A|/(N(1-\alpha))$.}{The outsider standing above which the
  platform admits (M15).}
  {No}{Computed from its inputs. In records it is
  bracketed by the outsider standings at which admissions start.}
\end{longtable}
\normalsize

\section*{Structure}

\begin{itemize}
\item \emph{Entering only through others.} The likelihood enters M2 and M3
  only through group sufficiency. $W$, $N$, $\alpha$, $|A|$ and the two
  standings enter M15 only through the admission margin.
\item \emph{Given directly by recorded comparisons.} The visible mean before
  and after members pass. The platform's admission decision.
\item \emph{Needing a computed reference.} The admission threshold $c_E$,
  from $W$, $N$, $\alpha$, $|A|$ and the host's standing by M15.
\item \emph{Elicited and not recorded.} Every belief, the prior, the group
  weights, the three posteriors and the salience. Each is the observer's and
  rationalises the observer's refusals, which are what records reach.
\item \emph{Assumed.} Group sufficiency, the quadratic payoff over the
  activity level and private consumption behind the engagement utility, the
  pair's level as the minimiser of the two losses, and the admission
  objective linear in the admission rate.
\item \emph{Never observed.} The passing gain, the passing cost, the
  baseline value and the outside option behind the engagement threshold,
  and the activity value and the two curvatures behind the sensitivity.
\end{itemize}

\end{document}
